\chapter*{Resumen}

El presente anteproyecto de graduación, titulado \textit{Diseño e implementación
de un sistema inteligente de bajo costo para el control ESD en laboratorios de
electrónica}, propone el desarrollo de una estación orientada a mejorar la
trazabilidad de las verificaciones ESD en laboratorios de electrónica. El
proyecto surge a partir de una oportunidad de mejora identificada en el
laboratorio TEM de Intel Costa Rica, donde actualmente se utiliza un probador
comercial de pulseras antiestáticas, pero no se dispone de un sistema integrado
que asocie cada prueba con un usuario, almacene registros y genere reportes
históricos para facilitar el seguimiento del procedimiento.

La solución propuesta integra un probador comercial de pulseras antiestáticas
como subsistema principal de verificación ESD, identificación de usuarios
mediante lectura de código de barras y reconocimiento facial complementario,
una interfaz gráfica táctil, almacenamiento local de registros y generación de
reportes históricos. Cada registro incluirá el nombre y el identificador del
usuario, la fecha, la hora y el resultado de la prueba ESD. Además, incorporará
la temperatura y la humedad relativa como información complementaria.

El sistema se plantea mediante una arquitectura modular y personalizable, con
un presupuesto de referencia aproximado de 300~USD, inferior al costo de las
soluciones comerciales avanzadas analizadas en el estado del arte. La
validación considerará los resultados de las pruebas funcionales individuales
de los subsistemas y la ejecución de pruebas de integración del sistema
completo. Estas pruebas están planificadas en el laboratorio TEM de Intel Costa
Rica. Si circunstancias fuera del control del proyecto impiden utilizar este
espacio, se aplicará el mismo procedimiento en un laboratorio de la Escuela de
Ingeniería Electrónica del Tecnológico de Costa Rica, Campus Tecnológico Local
San Carlos.

El proyecto no pretende sustituir sistemas industriales certificados, sino
demostrar la viabilidad de una estación ESD de bajo costo y adaptable para
mejorar el seguimiento del cumplimiento de los procedimientos ESD en
laboratorios de electrónica.

\textbf{Palabras clave:} ESD, control ESD, pulsera antiestática, trazabilidad,
laboratorio de electrónica, bajo costo, interfaz gráfica, identificación de
usuarios, reconocimiento facial, registros históricos.


\chapter*{Abstract}

This graduation project proposal, entitled \textit{Design and Implementation
of a Low-Cost Smart System for ESD Control in Electronics Laboratories},
proposes the development of a station aimed at improving the traceability of
ESD verifications in electronics laboratories. The project arises from an
improvement opportunity identified at the Intel Costa Rica TEM laboratory,
where a commercial wrist strap tester is currently used, but no integrated
system is available to associate each test with a user, store records, and
generate historical reports to facilitate monitoring of the procedure.

The proposed solution integrates a commercial wrist strap tester as the main
ESD verification subsystem, user identification through barcode reading and
complementary facial recognition, a touchscreen graphical interface, local
record storage, and historical report generation. Each record will include the
user's name and identifier, the date, time, and ESD test result. In addition,
temperature and relative humidity will be incorporated as supplementary
information.

The system is based on a modular and customizable architecture, with a
reference budget of approximately USD~300, lower than the cost of the advanced
commercial solutions analyzed in the state of the art. Validation will consider
the results of the individual functional tests of the subsystems and the
execution of complete-system integration tests. These tests are planned for the
Intel Costa Rica TEM laboratory. If circumstances beyond the project's control
prevent the use of this facility, the same procedure will be applied in a
laboratory of the School of Electronic Engineering at the Tecnológico de Costa
Rica, San Carlos Local Campus.

The project does not aim to replace certified industrial systems, but rather to
demonstrate the feasibility of a low-cost, adaptable ESD station for improving
the monitoring of ESD procedure compliance in electronics laboratories.

\textbf{Keywords:} ESD, ESD control, anti-static wrist strap, traceability,
electronics laboratory, low cost, graphical interface, user identification,
facial recognition, historical records.
